\begin{docpart}
    Calcul différentiel
\end{docpart}

\begin{exercise}[points=4]
    Soit $E$ un espace vectoriel normé et $B \colon E \times E \to \R$ une
    application bilinéaire continue. On pose $Q(x) = B(x,x)$.
    \begin{question}
        Montrer que $Q$ est différentiable et calculer $\Dif Q(x)$.
    \end{question}
    \begin{question}
        En déduire $\Dif^2 Q(x)$.
    \end{question}
\solution
    \begin{question}
        $Q(x+h) - Q(x) = B(x,h) + B(h,x) + B(h,h)$, et
        $\norm{B(h,h)} \le \norm{B}\, \norm{h}^2 = \smallo{h}$. Donc
        $\Dif Q(x) \cdot h = B(x,h) + B(h,x)$.
    \end{question}
    \begin{question}
        $\Dif Q$ est linéaire en $x$, donc
        $\Dif^2 Q(x)\cdot(h,k) = B(k,h) + B(h,k)$, indépendante de $x$.
    \end{question}
\end{exercise}

%---------------------------------------------------------------------------
\begin{docpart}
    Équations différentielles
\end{docpart}

\begin{exercise}[points=6]
    On considère $\dot{x}(t) = A\, x(t)$ avec
    $A = \begin{pmatrix} 0 & 1 \\ -1 & 0 \end{pmatrix}$.
    \begin{question}
        Calculer $e^{tA}$.
    \end{question}
    \begin{question}
        Décrire les orbites dans le plan de phase.
    \end{question}
\solution
    \begin{question}
        $A^2 = -I$, d'où
        $e^{tA} = \begin{pmatrix} \cos t & \sin t \\ -\sin t & \cos t\end{pmatrix}$.
    \end{question}
    \begin{question}
        Ce sont les cercles centrés à l'origine, parcourus dans le sens
        horaire à vitesse angulaire unité.
    \end{question}
\end{exercise}

\begin{exercise}[points=4, nosolution]
    Énoncer le théorème de Cauchy--Lipschitz et préciser ce que devient la
    conclusion si l'hypothèse de caractère localement lipschitzien est
    remplacée par la seule continuité.
\end{exercise}
